% =====================================================================
% Seção de Metodologia — versão genérica, RASCUNHO para o grupo.
% Incluída em proposal.tex por % =====================================================================
% Seção de Metodologia — versão genérica, RASCUNHO para o grupo.
% Incluída em proposal.tex por % =====================================================================
% Seção de Metodologia — versão genérica, RASCUNHO para o grupo.
% Incluída em proposal.tex por % =====================================================================
% Seção de Metodologia — versão genérica, RASCUNHO para o grupo.
% Incluída em proposal.tex por \input{metodologia}.
% =====================================================================

\section{Metodologia}

O objeto é avaliar um sistema web genérico, em que a requisição parte do
front-end, atravessa uma ou mais APIs e chega ao banco de dados. A análise
mede o serviço do ponto de vista do usuário e cada componente do fluxo,
para indicar em qual deles as falhas se concentram. Por meio desta análise
distribuída, objetiva-se desenvolver um plugin para a ferramenta Grafana.

\textbf{Coleta.} Os componentes são instrumentados com o OpenTelemetry,
cujas convenções semânticas padronizam as medidas de HTTP e de banco de
dados em qualquer sistema \cite{otel_site}. A métrica
\texttt{http.server.request.duration} registra código de resposta, rota e
duração de cada requisição. Os \textit{spans} de cliente de banco registram
as consultas feitas pela API. Quando o banco não é instrumentado, sua falha
é inferida pelo sinal que ela produz na API. Um \textit{plugin} do
Grafana desenvolvido pelo grupo calcula os índices por janela, e as
contagens exportadas em CSV para uma planilha servem para conferir os
valores do \textit{plugin}. A coleta usa tráfego real durante
aproximadamente 15 dias, condicionada à autorização do responsável por
cada fonte. Os dados são sempre mascarados: coletam-se apenas código de
resposta, rotas sem identificadores e as métricas de interesse.

\textbf{Critério de falha.} A falha é medida do ponto de vista do usuário,
como falha do serviço, e não só do componente \cite{oppenheimer2003}. Há
dois modos de falha. No modo \textit{erro}, a requisição recebe código 5xx
ou o código que a fonte usa para sinalizar falha percebida (por exemplo,
401 em massa quando a perda do banco encerra sessões). Desse código, conta
só o excesso sobre a taxa típica da rota, porque parte dele é uso normal
\cite{goseva2004}. No modo \textit{lentidão}, a duração passa de
$\tau = 200$~ms. O tempo é dividido em janelas de $\Delta = 5$~min, e só
entram janelas com pelo menos 100 requisições, para que uma única falha
não derrube o índice. Os SLOs são fixados antes da coleta, sem usar o
desempenho observado \cite{jones2016slo}: $R^{*}_{\mathrm{erro}} = 99\%$ e
$R^{*}_{\mathrm{lat}} = 95\%$.

\textbf{I1 — confiabilidade por requisição.} Pelo modelo de Nelson, usado
por \citeonline{tian2004} e \citeonline{goseva2004}, $R = 1-f/n$, em que $f$ é o número de requisições com falha e $n$ o total de
requisições. $R$ é a probabilidade de uma requisição ser atendida
corretamente. A razão normaliza as falhas pela carga \cite{tian2004} e
aproxima a indisponibilidade percebida pelo usuário
\cite{alvidrez2016risk}. O índice é calculado para cada modo de falha
($R_{\mathrm{erro}}$ e $R_{\mathrm{lat}}$) e para os dois juntos, no
período inteiro, por janela ($R_j$) e por componente ($R_i$). Os percentis
95 e 99 da latência complementam $R_{\mathrm{lat}}$ como medida de
variabilidade.

\textbf{I2 — disponibilidade em relação ao SLO.} A janela $j$ está em
falha quando $R_{\mathrm{erro},j} < R^{*}_{\mathrm{erro}}$ ou
$R_{\mathrm{lat},j} < R^{*}_{\mathrm{lat}}$. A disponibilidade $A$ é a fração de janelas
dentro do SLO.

\textbf{I3 — MTBF e MTTR.} Uma sequência de janelas consecutivas em falha
forma um episódio de falha. O MTBF é o tempo total dentro do SLO dividido
pelo número de episódios, e o MTTR é a duração média dos episódios
\cite{oconnor2012}. Como $A = \mathrm{MTBF}/(\mathrm{MTBF}+\mathrm{MTTR})$, os dois
mostram se o serviço perde disponibilidade por falhar com frequência ou
por demorar a se recuperar.

\textbf{Declaração de uso de IA.}
% [AJUSTAR] Completem com todas as ferramentas que o grupo usou.
Ferramenta: Claude Opus 5.5 (Anthropic), em novembro de 2026. Finalidade:
busca de referências, discussão de índices e rascunho desta seção.
Estratégia: o grupo revisou e reescreveu o texto e conferiu as referências
nas fontes originais. Responsabilidade: o conteúdo é de responsabilidade
exclusiva dos autores.
.
% =====================================================================

\section{Metodologia}

O objeto é avaliar um sistema web genérico, em que a requisição parte do
front-end, atravessa uma ou mais APIs e chega ao banco de dados. A análise
mede o serviço do ponto de vista do usuário e cada componente do fluxo,
para indicar em qual deles as falhas se concentram. Por meio desta análise
distribuída, objetiva-se desenvolver um plugin para a ferramenta Grafana.

\textbf{Coleta.} Os componentes são instrumentados com o OpenTelemetry,
cujas convenções semânticas padronizam as medidas de HTTP e de banco de
dados em qualquer sistema \cite{otel_site}. A métrica
\texttt{http.server.request.duration} registra código de resposta, rota e
duração de cada requisição. Os \textit{spans} de cliente de banco registram
as consultas feitas pela API. Quando o banco não é instrumentado, sua falha
é inferida pelo sinal que ela produz na API. Um \textit{plugin} do
Grafana desenvolvido pelo grupo calcula os índices por janela, e as
contagens exportadas em CSV para uma planilha servem para conferir os
valores do \textit{plugin}. A coleta usa tráfego real durante
aproximadamente 15 dias, condicionada à autorização do responsável por
cada fonte. Os dados são sempre mascarados: coletam-se apenas código de
resposta, rotas sem identificadores e as métricas de interesse.

\textbf{Critério de falha.} A falha é medida do ponto de vista do usuário,
como falha do serviço, e não só do componente \cite{oppenheimer2003}. Há
dois modos de falha. No modo \textit{erro}, a requisição recebe código 5xx
ou o código que a fonte usa para sinalizar falha percebida (por exemplo,
401 em massa quando a perda do banco encerra sessões). Desse código, conta
só o excesso sobre a taxa típica da rota, porque parte dele é uso normal
\cite{goseva2004}. No modo \textit{lentidão}, a duração passa de
$\tau = 200$~ms. O tempo é dividido em janelas de $\Delta = 5$~min, e só
entram janelas com pelo menos 100 requisições, para que uma única falha
não derrube o índice. Os SLOs são fixados antes da coleta, sem usar o
desempenho observado \cite{jones2016slo}: $R^{*}_{\mathrm{erro}} = 99\%$ e
$R^{*}_{\mathrm{lat}} = 95\%$.

\textbf{I1 — confiabilidade por requisição.} Pelo modelo de Nelson, usado
por \citeonline{tian2004} e \citeonline{goseva2004}, $R = 1-f/n$, em que $f$ é o número de requisições com falha e $n$ o total de
requisições. $R$ é a probabilidade de uma requisição ser atendida
corretamente. A razão normaliza as falhas pela carga \cite{tian2004} e
aproxima a indisponibilidade percebida pelo usuário
\cite{alvidrez2016risk}. O índice é calculado para cada modo de falha
($R_{\mathrm{erro}}$ e $R_{\mathrm{lat}}$) e para os dois juntos, no
período inteiro, por janela ($R_j$) e por componente ($R_i$). Os percentis
95 e 99 da latência complementam $R_{\mathrm{lat}}$ como medida de
variabilidade.

\textbf{I2 — disponibilidade em relação ao SLO.} A janela $j$ está em
falha quando $R_{\mathrm{erro},j} < R^{*}_{\mathrm{erro}}$ ou
$R_{\mathrm{lat},j} < R^{*}_{\mathrm{lat}}$. A disponibilidade $A$ é a fração de janelas
dentro do SLO.

\textbf{I3 — MTBF e MTTR.} Uma sequência de janelas consecutivas em falha
forma um episódio de falha. O MTBF é o tempo total dentro do SLO dividido
pelo número de episódios, e o MTTR é a duração média dos episódios
\cite{oconnor2012}. Como $A = \mathrm{MTBF}/(\mathrm{MTBF}+\mathrm{MTTR})$, os dois
mostram se o serviço perde disponibilidade por falhar com frequência ou
por demorar a se recuperar.

\textbf{Declaração de uso de IA.}
% [AJUSTAR] Completem com todas as ferramentas que o grupo usou.
Ferramenta: Claude Opus 5.5 (Anthropic), em novembro de 2026. Finalidade:
busca de referências, discussão de índices e rascunho desta seção.
Estratégia: o grupo revisou e reescreveu o texto e conferiu as referências
nas fontes originais. Responsabilidade: o conteúdo é de responsabilidade
exclusiva dos autores.
.
% =====================================================================

\section{Metodologia}

O objeto é avaliar um sistema web genérico, em que a requisição parte do
front-end, atravessa uma ou mais APIs e chega ao banco de dados. A análise
mede o serviço do ponto de vista do usuário e cada componente do fluxo,
para indicar em qual deles as falhas se concentram. Por meio desta análise
distribuída, objetiva-se desenvolver um plugin para a ferramenta Grafana.

\textbf{Coleta.} Os componentes são instrumentados com o OpenTelemetry,
cujas convenções semânticas padronizam as medidas de HTTP e de banco de
dados em qualquer sistema \cite{otel_site}. A métrica
\texttt{http.server.request.duration} registra código de resposta, rota e
duração de cada requisição. Os \textit{spans} de cliente de banco registram
as consultas feitas pela API. Quando o banco não é instrumentado, sua falha
é inferida pelo sinal que ela produz na API. Um \textit{plugin} do
Grafana desenvolvido pelo grupo calcula os índices por janela, e as
contagens exportadas em CSV para uma planilha servem para conferir os
valores do \textit{plugin}. A coleta usa tráfego real durante
aproximadamente 15 dias, condicionada à autorização do responsável por
cada fonte. Os dados são sempre mascarados: coletam-se apenas código de
resposta, rotas sem identificadores e as métricas de interesse.

\textbf{Critério de falha.} A falha é medida do ponto de vista do usuário,
como falha do serviço, e não só do componente \cite{oppenheimer2003}. Há
dois modos de falha. No modo \textit{erro}, a requisição recebe código 5xx
ou o código que a fonte usa para sinalizar falha percebida (por exemplo,
401 em massa quando a perda do banco encerra sessões). Desse código, conta
só o excesso sobre a taxa típica da rota, porque parte dele é uso normal
\cite{goseva2004}. No modo \textit{lentidão}, a duração passa de
$\tau = 200$~ms. O tempo é dividido em janelas de $\Delta = 5$~min, e só
entram janelas com pelo menos 100 requisições, para que uma única falha
não derrube o índice. Os SLOs são fixados antes da coleta, sem usar o
desempenho observado \cite{jones2016slo}: $R^{*}_{\mathrm{erro}} = 99\%$ e
$R^{*}_{\mathrm{lat}} = 95\%$.

\textbf{I1 — confiabilidade por requisição.} Pelo modelo de Nelson, usado
por \citeonline{tian2004} e \citeonline{goseva2004}, $R = 1-f/n$, em que $f$ é o número de requisições com falha e $n$ o total de
requisições. $R$ é a probabilidade de uma requisição ser atendida
corretamente. A razão normaliza as falhas pela carga \cite{tian2004} e
aproxima a indisponibilidade percebida pelo usuário
\cite{alvidrez2016risk}. O índice é calculado para cada modo de falha
($R_{\mathrm{erro}}$ e $R_{\mathrm{lat}}$) e para os dois juntos, no
período inteiro, por janela ($R_j$) e por componente ($R_i$). Os percentis
95 e 99 da latência complementam $R_{\mathrm{lat}}$ como medida de
variabilidade.

\textbf{I2 — disponibilidade em relação ao SLO.} A janela $j$ está em
falha quando $R_{\mathrm{erro},j} < R^{*}_{\mathrm{erro}}$ ou
$R_{\mathrm{lat},j} < R^{*}_{\mathrm{lat}}$. A disponibilidade $A$ é a fração de janelas
dentro do SLO.

\textbf{I3 — MTBF e MTTR.} Uma sequência de janelas consecutivas em falha
forma um episódio de falha. O MTBF é o tempo total dentro do SLO dividido
pelo número de episódios, e o MTTR é a duração média dos episódios
\cite{oconnor2012}. Como $A = \mathrm{MTBF}/(\mathrm{MTBF}+\mathrm{MTTR})$, os dois
mostram se o serviço perde disponibilidade por falhar com frequência ou
por demorar a se recuperar.

\textbf{Declaração de uso de IA.}
% [AJUSTAR] Completem com todas as ferramentas que o grupo usou.
Ferramenta: Claude Opus 5.5 (Anthropic), em novembro de 2026. Finalidade:
busca de referências, discussão de índices e rascunho desta seção.
Estratégia: o grupo revisou e reescreveu o texto e conferiu as referências
nas fontes originais. Responsabilidade: o conteúdo é de responsabilidade
exclusiva dos autores.
.
% =====================================================================

\section{Metodologia}

O objeto é avaliar um sistema web genérico, em que a requisição parte do
front-end, atravessa uma ou mais APIs e chega ao banco de dados. A análise
mede o serviço do ponto de vista do usuário e cada componente do fluxo,
para indicar em qual deles as falhas se concentram. Por meio desta análise
distribuída, objetiva-se desenvolver um plugin para a ferramenta Grafana.

\textbf{Coleta.} Os componentes são instrumentados com o OpenTelemetry,
cujas convenções semânticas padronizam as medidas de HTTP e de banco de
dados em qualquer sistema \cite{otel_site}. A métrica
\texttt{http.server.request.duration} registra código de resposta, rota e
duração de cada requisição. Os \textit{spans} de cliente de banco registram
as consultas feitas pela API. Quando o banco não é instrumentado, sua falha
é inferida pelo sinal que ela produz na API. Um \textit{plugin} do
Grafana desenvolvido pelo grupo calcula os índices por janela, e as
contagens exportadas em CSV para uma planilha servem para conferir os
valores do \textit{plugin}. A coleta usa tráfego real durante
aproximadamente 15 dias, condicionada à autorização do responsável por
cada fonte. Os dados são sempre mascarados: coletam-se apenas código de
resposta, rotas sem identificadores e as métricas de interesse.

\textbf{Critério de falha.} A falha é medida do ponto de vista do usuário,
como falha do serviço, e não só do componente \cite{oppenheimer2003}. Há
dois modos de falha. No modo \textit{erro}, a requisição recebe código 5xx
ou o código que a fonte usa para sinalizar falha percebida (por exemplo,
401 em massa quando a perda do banco encerra sessões). Desse código, conta
só o excesso sobre a taxa típica da rota, porque parte dele é uso normal
\cite{goseva2004}. No modo \textit{lentidão}, a duração passa de
$\tau = 200$~ms. O tempo é dividido em janelas de $\Delta = 5$~min, e só
entram janelas com pelo menos 100 requisições, para que uma única falha
não derrube o índice. Os SLOs são fixados antes da coleta, sem usar o
desempenho observado \cite{jones2016slo}: $R^{*}_{\mathrm{erro}} = 99\%$ e
$R^{*}_{\mathrm{lat}} = 95\%$.

\textbf{I1 — confiabilidade por requisição.} Pelo modelo de Nelson, usado
por \citeonline{tian2004} e \citeonline{goseva2004}, $R = 1-f/n$, em que $f$ é o número de requisições com falha e $n$ o total de
requisições. $R$ é a probabilidade de uma requisição ser atendida
corretamente. A razão normaliza as falhas pela carga \cite{tian2004} e
aproxima a indisponibilidade percebida pelo usuário
\cite{alvidrez2016risk}. O índice é calculado para cada modo de falha
($R_{\mathrm{erro}}$ e $R_{\mathrm{lat}}$) e para os dois juntos, no
período inteiro, por janela ($R_j$) e por componente ($R_i$). Os percentis
95 e 99 da latência complementam $R_{\mathrm{lat}}$ como medida de
variabilidade.

\textbf{I2 — disponibilidade em relação ao SLO.} A janela $j$ está em
falha quando $R_{\mathrm{erro},j} < R^{*}_{\mathrm{erro}}$ ou
$R_{\mathrm{lat},j} < R^{*}_{\mathrm{lat}}$. A disponibilidade $A$ é a fração de janelas
dentro do SLO.

\textbf{I3 — MTBF e MTTR.} Uma sequência de janelas consecutivas em falha
forma um episódio de falha. O MTBF é o tempo total dentro do SLO dividido
pelo número de episódios, e o MTTR é a duração média dos episódios
\cite{oconnor2012}. Como $A = \mathrm{MTBF}/(\mathrm{MTBF}+\mathrm{MTTR})$, os dois
mostram se o serviço perde disponibilidade por falhar com frequência ou
por demorar a se recuperar.

\textbf{Declaração de uso de IA.}
% [AJUSTAR] Completem com todas as ferramentas que o grupo usou.
Ferramenta: Claude Opus 5.5 (Anthropic), em novembro de 2026. Finalidade:
busca de referências, discussão de índices e rascunho desta seção.
Estratégia: o grupo revisou e reescreveu o texto e conferiu as referências
nas fontes originais. Responsabilidade: o conteúdo é de responsabilidade
exclusiva dos autores.
